% Jeans-Independent Holdback Bypass for AMR Refinement
% Technical note — cuRAMSES / Horizon-5
%
\documentclass[11pt,a4paper]{article}

\usepackage{amsmath,amssymb}
\usepackage{graphicx}
\usepackage{booktabs}
\usepackage{hyperref}
\usepackage{xcolor}
\usepackage{listings}
\usepackage{natbib}

\lstset{
  language=[90]Fortran,
  basicstyle=\ttfamily\small,
  keywordstyle=\color{blue},
  commentstyle=\color{gray},
  frame=single,
  breaklines=true
}

\newcommand{\lJ}{\lambda_{\rm J}}
\newcommand{\nstar}{n_\star}
\newcommand{\Tstar}{T_{2,\star}}
\newcommand{\gstar}{g_\star}
\newcommand{\dx}{\Delta x}
\newcommand{\ramses}{\textsc{ramses}}
\newcommand{\curamses}{\textsc{cuRAMSES}}
\newcommand{\nlevelmax}{N_{\rm levelmax}}
\newcommand{\twotondim}{2^{N_{\rm dim}}}
\newcommand{\code}[1]{\texttt{#1}}

\title{Jeans-Independent Holdback Bypass\\for AMR Refinement in Cosmological Simulations}
\author{Juhan Kim (KIAS)}
\date{March 2026}

\begin{document}
\maketitle

\begin{abstract}
Cosmological AMR simulations employ a \emph{holdback} mechanism that
progressively releases refinement levels as the universe expands,
preventing premature over-refinement at high redshift.  However, this
mechanism blocks \emph{all} refinement criteria simultaneously, including
the Jeans length criterion, which is a numerical stability requirement
rather than merely a resolution preference.  When a holdback level is
finally released, a large number of cells suddenly satisfy the
refinement condition, triggering an abrupt starburst.  We implement a
\code{jeans\_bypass\_holdback} option that allows the Jeans criterion to
operate independently of the holdback, with a configurable level limit
(\code{n\_jeans\_bypass\_levels}) to prevent uncontrolled refinement
cascades.  We discuss the interplay between the Jeans criterion, the
polytropic equation of state, and the density threshold
\code{d\_jeans\_thre}, showing that the polytropic temperature floor
provides a natural self-regulation mechanism for the bypass.
\end{abstract}

%======================================================================
\section{AMR Refinement in \ramses}
\label{sec:amr_refine}
%======================================================================

In \ramses, each cell at level $\ell < \nlevelmax$ is tested for
refinement at every coarse time-step.  The decision to subdivide a
cell into $\twotondim$ children is governed by a combination of
physical, geometric, and cosmological criteria.

\subsection{Refinement logic}
\label{sec:refine_logic}

The overall refinement flag for a cell is determined by the following
Boolean expression:
\begin{equation}
\mathcal{R} = \lnot\,\mathcal{H}
  \;\land\; \mathcal{G}
  \;\land\; \bigl(
    \mathcal{M} \lor \mathcal{S} \lor \mathcal{D} \lor \mathcal{P}
    \lor \mathcal{U} \lor \mathcal{J}
  \bigr),
\label{eq:refine_flag}
\end{equation}
where each symbol is defined below.

\subsubsection{Physical criteria (OR)}

\begin{enumerate}

\item \textbf{Quasi-Lagrangian mass criterion ($\mathcal{M}$).}
  The cell is flagged if its gas density exceeds a
  level-dependent threshold:
  \begin{equation}
    \rho_{\rm cell} \;\ge\; m_{\rm refine}(\ell)
      \;\frac{m_{\rm sph}}{V_{\rm cell}},
  \end{equation}
  where $m_{\rm sph}$ is the effective SPH kernel mass and
  $V_{\rm cell} = (\dx_\ell)^3$.
  This criterion ensures that regions containing roughly
  $m_{\rm refine}$ or more particles per cell are refined,
  providing a quasi-Lagrangian resolution that tracks gravitational
  collapse.  Controlled by the namelist array
  \code{m\_refine(1:nlevelmax)}.

\item \textbf{Sink particle forcing ($\mathcal{S}$).}
  If \code{sink\_refine=.true.}, any cell containing a sink cloud
  particle (identified by $\code{idp} < 0$ and $\code{tp} = 0$) is
  unconditionally flagged up to $\nlevelmax$, guaranteeing maximum
  resolution around each accretion region.

\item \textbf{Density gradient ($\mathcal{D}$).}
  The normalised density gradient between a cell and each of its
  $2N_{\rm dim}$ face-neighbours is computed as
  \begin{equation}
    \epsilon_\rho = 2\,\max\!\left(
      \frac{|\rho_+ - \rho_0|}{|\rho_+ + \rho_0| + \rho_{\rm floor}},\;
      \frac{|\rho_0 - \rho_-|}{|\rho_0 + \rho_-| + \rho_{\rm floor}}
    \right).
  \end{equation}
  If $\epsilon_\rho > \code{err\_grad\_d}$ along any axis, the cell
  is flagged.  This captures density discontinuities such as shocks
  and contact waves.

\item \textbf{Pressure gradient ($\mathcal{P}$).}
  Analogous to the density gradient but applied to the thermal
  pressure:
  \begin{equation}
    \epsilon_P = 2\,\max\!\left(
      \frac{|P_+ - P_0|}{|P_+ + P_0| + P_{\rm floor}},\;
      \frac{|P_0 - P_-|}{|P_0 + P_-| + P_{\rm floor}}
    \right),
  \end{equation}
  flagging when $\epsilon_P > \code{err\_grad\_p}$.

\item \textbf{Velocity gradient ($\mathcal{U}$).}
  A normalised velocity jump relative to the local sound speed:
  \begin{equation}
    \epsilon_v = 2\,\max\!\left(
      \frac{|v_+ - v_0|}{c_{s,+} + c_{s,0} + |v_+| + |v_0|
                          + v_{\rm floor}},\;
      \frac{|v_0 - v_-|}{c_{s,0} + c_{s,-} + |v_0| + |v_-|
                          + v_{\rm floor}}
    \right),
  \end{equation}
  flagging when $\epsilon_v > \code{err\_grad\_u}$ in any component.
  This criterion targets strong velocity shears and accretion shocks.

\item \textbf{Jeans length ($\mathcal{J}$).}
  When \code{jeans\_refine($\ell$)} $= N_J > 0$, the cell is flagged if
  the local Jeans length is shorter than $N_J$ cell widths:
  \begin{equation}
    \lJ = \sqrt{\frac{\pi\,c_s^2}{G\,\rho}}
    < N_J\,\dx_\ell,
    \label{eq:jeans_crit}
  \end{equation}
  implementing the \citet{Truelove1997} condition to prevent artificial
  fragmentation in self-gravitating gas.

  An additional density threshold \code{d\_jeans\_thre} restricts this
  criterion to cells with $n_{\rm H} \ge \code{d\_jeans\_thre}$.

\end{enumerate}

\subsubsection{Geometry mask (AND)}

\textbf{Ellipsoidal region ($\mathcal{G}$).}
When \code{r\_refine($\ell$)} $> -1$, only cells inside a generalised
ellipsoidal region are permitted to refine.  The test is
\begin{equation}
  \left(
    \frac{2|\Delta x|}{r}
  \right)^{e_r}
  + \left(
    \frac{2|\Delta y|}{a\,r}
  \right)^{e_r}
  + \left(
    \frac{2|\Delta z|}{b\,r}
  \right)^{e_r}
  < 1,
\end{equation}
where $r = \code{r\_refine}$ is the region diameter,
$a = \code{a\_refine}$ and $b = \code{b\_refine}$ are the
axis ratios, and $e_r = \code{exp\_refine}$ is the norm exponent
($e_r = 2$ gives an ellipsoid; $e_r \to \infty$ a rectangular box).
Cells outside this region are \emph{un-flagged} regardless of
physical criteria, restricting high resolution to a zoom-in target.

\subsubsection{Cosmological holdback (NOT)}

\textbf{Resolution holdback ($\mathcal{H}$).}
In cosmological runs with cooling (\code{cosmo=.true.},
\code{cooling=.true.}), an overly eager refinement at high redshift
can exhaust memory before the epoch of interest.
When \code{q\_refine\_holdback=.true.} and the physical cell size
$\dx_\ell < 4^{1/N_{\rm dim}}\,\dx_{\nlevelmax}/a_{\rm exp}$
for levels $\ell > \code{nlevelmax\_part} + 3$, the \emph{entire}
level is blocked from further refinement:
\begin{equation}
  \mathcal{H}(\ell) = \bigl[\ell > \ell_{\rm part} + \ell_{\rm collapse}\bigr]
  \;\wedge\;
  \bigl[\dx_\ell < 4^{1/N_{\rm dim}}\,\dx_{\rm min}/a\bigr],
  \label{eq:holdback}
\end{equation}
where $\ell_{\rm part}$ is the particle initial-condition level
(typically $\ell_{\rm min}$) and $\ell_{\rm collapse}=3$.
This enforces a constant physical resolution at low redshift and a
constant comoving resolution at high redshift, gradually releasing
deeper levels as the universe expands.

The first holdback level $\ell_{\rm hb}$ is the smallest $\ell$ satisfying
Eq.~\eqref{eq:holdback}:
\begin{equation}
  \ell_{\rm hb} = \min\bigl\{\ell > \ell_{\rm part}+\ell_{\rm collapse}
  \;:\; \dx_\ell < 4^{1/N_{\rm dim}}\,\dx_{\rm min}/a\bigr\}.
  \label{eq:lhb}
\end{equation}

\subsection{Buffer zones}

After applying equation~(\ref{eq:refine_flag}), flagged cells are
expanded by \code{nexpand($\ell$)} cell widths in each direction via a
dilation operator (\code{smooth\_fine}).
This cubic buffer ensures smooth refinement boundaries, preventing
the Godunov solver from encountering abrupt level transitions that
would degrade numerical accuracy.


%======================================================================
\section{The Starburst Problem}
\label{sec:starburst}
%======================================================================

The holdback mechanism blocks \emph{all} refinement criteria at levels
$\ell \ge \ell_{\rm hb}$, including the Jeans criterion.  This creates
two related problems:

\paragraph{Artificial fragmentation.}
The Jeans length criterion $\mathcal{J}$ differs fundamentally from
the other refinement triggers: it is a \emph{numerical stability}
requirement \citep{Truelove1997}, not merely a resolution preference.
Violating the Jeans condition leads to artificial fragmentation ---
a qualitatively wrong result --- regardless of epoch.  Yet in the
standard \ramses\ logic (equation~\ref{eq:refine_flag}), the
cosmological holdback $\mathcal{H}$ blocks $\mathcal{J}$ uniformly
with all other criteria.  At high redshift, dense self-gravitating
regions may therefore be held at insufficient resolution, producing
spurious fragmentation that propagates into lower-redshift structure.

\paragraph{Abrupt SFR transition.}
When the holdback releases level $\ell$ (i.e.\ $\mathcal{H}(\ell)$
transitions from true to false), all accumulated criteria fire
simultaneously.  Many cells that have been gravitationally collapsing
during the holdback period now suddenly satisfy the refinement
condition, leading to a burst of new cells and an abrupt spike in
the star formation rate.


%======================================================================
\section{Jeans Bypass: Design and Implementation}
\label{sec:bypass}
%======================================================================

\subsection{Modified refinement logic}

The key observation is that the Jeans length criterion is
\emph{cell-local}: it depends only on the density and temperature of the
cell itself (\code{uold}), with no neighbor information required.  This
allows a lightweight evaluation even when the full refinement machinery
(which requires neighbor gathering for $\mathcal{D}$, $\mathcal{P}$,
$\mathcal{U}$) is skipped by the holdback.

With \code{jeans\_bypass\_holdback=.true.}, the refinement logic becomes
\begin{equation}
\mathcal{R}' = \Bigl[
  \lnot\,\mathcal{H} \;\land\; \mathcal{G} \;\land\;
  \bigl(\mathcal{M} \lor \mathcal{S} \lor \mathcal{D}
        \lor \mathcal{P} \lor \mathcal{U} \lor \mathcal{J}\bigr)
\Bigr]
\;\lor\;
\Bigl[ \mathcal{H} \;\land\; \mathcal{J} \;\land\;
  (\ell < \ell_{\rm hb} + n_{\rm bypass}) \Bigr],
\label{eq:refine_bypass}
\end{equation}
where $n_{\rm bypass} \equiv$ \code{n\_jeans\_bypass\_levels} (default 1).
The geometry mask $\mathcal{G}$ still applies through the non-holdback
branch, preserving zoom-in spatial restrictions.  Only the temporal
holdback $\mathcal{H}$ is bypassed for $\mathcal{J}$, and only within
a limited number of levels.

This modification is analogous to the existing treatment of sink
particles: $\mathcal{S}$ is already evaluated inside the
holdback-protected block in the current code, but unlike $\mathcal{J}$
it only triggers in cells that already host a sink.  Extending the same
exemption to $\mathcal{J}$ ensures numerical integrity for gas that is
\emph{about to} form a sink.


\subsection{Level limitation: preventing cascades}
\label{sec:cascade}

Without a level limit ($n_{\rm bypass}\to\infty$), the Jeans bypass
creates a catastrophic cascade.  When a cell at the holdback level
$\ell_{\rm hb}$ is flagged by the Jeans criterion, child cells are
created at level $\ell_{\rm hb}+1$.  Since level $\ell_{\rm hb}+1$ is
also in the holdback zone, the bypass applies again, flagging cells at
$\ell_{\rm hb}+2$, and so on until $\ell_{\rm max}$.  This rapidly
exhausts available memory.

The cascade occurs because the Jeans length in radiatively cooling gas
(below the star-formation density threshold) shrinks without bound as
the temperature drops.  There is no physical mechanism to halt the
Jeans criterion from triggering at every successive level --- only the
polytropic temperature floor (Section~\ref{sec:polytropic}) or a hard
level limit can stop it.

Table~\ref{tab:cascade} shows the progression observed in a test run
without level limit:

\begin{table}[h]
\centering
\caption{Cascade refinement without level limit ($a=0.164$,
$\ell_{\rm hb}=14$, $\ell_{\rm max}=17$).  All levels $\ge 14$ are
in the holdback zone.}
\label{tab:cascade}
\begin{tabular}{cccccl}
\toprule
Level & $\dx/\dx_{\rm min}$ & Holdback & Bypass & $\mathcal{J}$ fires & Result \\
\midrule
$\le 12$ & $\ge 32$  & No  & --- & Yes (normal) & Normal \\
13       & 16        & No  & --- & Yes (normal) & Normal \\
14       & 8         & Yes & Active & Yes & $\to$ lv.~15 created \\
15       & 4         & Yes & Active & Yes & $\to$ lv.~16 created \\
16       & 2         & Yes & Active & Yes & $\to$ lv.~17 created \\
17       & 1         & $\ell_{\rm max}$ & --- & --- & OOM crash \\
\bottomrule
\end{tabular}
\end{table}

The memory footprint grew from 81\% to 88.5\% in a single coarse step
before crashing (Job 122644).  The fine-step timestep also decreased by
$8\times$ due to subcycling at levels 15--17.

With $n_{\rm bypass}=1$, only level $\ell_{\rm hb}=14$ is bypassed,
creating cells at level 15 but no further.  This provides a controlled
pre-refinement before the holdback releases.


%======================================================================
\section{Role of the Polytropic Temperature Floor}
\label{sec:polytropic}
%======================================================================

\subsection{Jeans length in the polytropic regime}

In \ramses, the star-formation module enforces a polytropic equation
of state above the density threshold $\nstar$:
\begin{equation}
  T(n) = \Tstar\,\Bigl(\frac{n}{\nstar}\Bigr)^{\gstar - 1},
  \qquad n \ge \nstar,
  \label{eq:polytropic}
\end{equation}
with default parameters $\Tstar = 750\,{\rm K}$,
$\nstar = 0.1\,{\rm cm}^{-3}$, $\gstar = 5/3$.  This floor is
applied in the cooling routine and is already reflected in \code{uold},
so the Jeans length computation in \code{jeans\_length\_refine()}
automatically includes the polytropic stabilisation.

Substituting Eq.~\eqref{eq:polytropic} into Eq.~\eqref{eq:jeans_crit}:
\begin{equation}
  \lJ \propto \sqrt{\frac{T}{\rho}} \propto
  \sqrt{\rho^{\gstar-1}/\rho} = \rho^{(\gstar-2)/2}.
  \label{eq:jeans_poly}
\end{equation}
For $\gstar = 5/3$: $\lJ \propto \rho^{-1/6}$.  The Jeans length
\emph{decreases very slowly} with density --- one or two additional
levels of refinement are typically sufficient to resolve it.

In contrast, for gas below $\nstar$ (no polytropic floor), radiative
cooling can drive $T$ to very low values while $\rho$ remains moderate.
This yields $\lJ \ll \dx_\ell$ at many levels, causing the cascade.

\subsection{Self-regulation via \code{d\_jeans\_thre}}

The density threshold parameter \code{d\_jeans\_thre} (in units of
$n_{\rm H}\,[{\rm cm}^{-3}]$) restricts the Jeans criterion to cells
above a given density:
\begin{equation}
  \mathcal{J}(\text{cell}) =
  \begin{cases}
    N_J\,\dx_\ell \ge \lJ & \text{if } n \ge
    \code{d\_jeans\_thre}, \\
    \text{false} & \text{otherwise}.
  \end{cases}
\end{equation}

Setting $\code{d\_jeans\_thre} \approx \nstar$ restricts the Jeans
criterion to the polytropic regime.  The cascade then self-limits
through the following mechanism:

\begin{enumerate}
  \item At level $\ell_{\rm hb}$: only cells with $n \ge \nstar$ are
    checked.  The polytropic floor guarantees $T \ge \Tstar$, so $\lJ$
    is relatively large.
  \item If $N_J\,\dx_{\ell_{\rm hb}} > \lJ$, the cell is flagged and
    children at $\ell_{\rm hb}+1$ are created.
  \item At level $\ell_{\rm hb}+1$: $\dx$ is halved, density is
    approximately unchanged (refinement preserves cell-averaged
    quantities), $\lJ$ is approximately unchanged.
  \item $N_J\,\dx_{\ell_{\rm hb}+1} = N_J\,\dx_{\ell_{\rm hb}}/2
    < \lJ$ $\Rightarrow$ Jeans criterion is \emph{satisfied}
    (the Jeans length is now resolved) $\Rightarrow$ no further
    flagging $\Rightarrow$ cascade stops naturally.
\end{enumerate}

This provides a \emph{physics-based} self-regulation, complementary to
the \emph{hard limit} imposed by \code{n\_jeans\_bypass\_levels}.

The two regulators are compared in Table~\ref{tab:regulators}.

\begin{table}[h]
\centering
\caption{Two complementary mechanisms for preventing the bypass cascade.}
\label{tab:regulators}
\begin{tabular}{lll}
\toprule
Mechanism & Type & Behaviour \\
\midrule
\code{n\_jeans\_bypass\_levels} & Hard limit &
  Caps extra levels regardless of physics \\
\code{d\_jeans\_thre} $\approx \nstar$ & Physics-based &
  Polytropic floor limits $\lJ$ decrease \\
\bottomrule
\end{tabular}
\end{table}

When \code{d\_jeans\_thre}$\,=0$ (default), \emph{all} cells are
checked, including low-density gas where cooling produces arbitrarily
small $\lJ$.  This is the root cause of the unlimited cascade observed
in Section~\ref{sec:cascade}.


%======================================================================
\section{Implementation Details}
\label{sec:implementation}
%======================================================================

\subsection{Code structure}

The implementation modifies two source files (plus parameter declaration
and namelist reading).

\paragraph{\code{flag\_utils.kjhan.f90} --- \code{userflag\_fine()}}
The holdback early-return is modified from a simple \code{return} to a
conditional block:

\begin{lstlisting}
if(prevent_refine .and. q_refine_holdback) then
   if(jeans_bypass_holdback .and. hydro &
      .and. n_jeans_bypass_levels > 0) then
      ! Find the first holdback level
      dx_min = (0.5D0**nlevelmax)*scale
      holdback_threshold = (4d0**(1d0/ndim))*(dx_min/aexp)
      holdback_first = nlevelmax + 1
      do jlevel = nlevelmax_part+nlevel_collapse+1, nlevelmax
         dx_test = (0.5d0**jlevel)*scale
         if(dx_test < holdback_threshold) then
            holdback_first = jlevel
            exit
         endif
      end do
      ! Only bypass within allowed range
      if(ilevel < holdback_first &
         + n_jeans_bypass_levels) then
         call jeans_only_flag(ilevel)
      endif
   endif
   return
endif
\end{lstlisting}

\paragraph{\code{hydro\_flag.kjhan.f90} --- \code{jeans\_only\_flag()}}
A new lightweight subroutine that:
\begin{enumerate}
  \item Loops over active grids with OpenMP parallelism.
  \item Computes cell indices without neighbor gathering.
  \item Calls the existing \code{jeans\_length\_refine()} subroutine
    (no code duplication).
  \item Sets \code{flag1} for flagged cells and updates \code{nflag}.
\end{enumerate}

This is possible because \code{jeans\_length\_refine()} only accesses
cell-local data (\code{uold(ind\_cell,...)}) and requires no neighbor
information, unlike the gradient-based criteria
$\mathcal{D}$, $\mathcal{P}$, $\mathcal{U}$.

\subsection{Namelist parameters}

In \code{amr\_parameters.jaehyun.f90}:
\begin{lstlisting}
logical :: jeans_bypass_holdback = .false.
integer :: n_jeans_bypass_levels = 1
\end{lstlisting}

Both are read via \code{\&REFINE\_PARAMS} in
\code{read\_hydro\_params.f90}.

\subsection{Safety properties}

\begin{itemize}
  \item \textbf{Default off}: With \code{jeans\_bypass\_holdback=.false.}
    (default), the new code path is never entered.  Results are
    bit-identical to the unmodified code.
  \item \textbf{No neighbor access}: \code{jeans\_only\_flag()} is
    safe to call at any point in the refinement cycle.
  \item \textbf{OMP safety}: \code{flag1} writes are grid-partitioned
    (no race conditions); \code{nflag} uses reduction.
  \item \textbf{Hard level limit}: \code{n\_jeans\_bypass\_levels}
    provides a fail-safe even if the physics-based regulation
    (\code{d\_jeans\_thre}) is misconfigured.
\end{itemize}


%======================================================================
\section{Summary of Parameters}
\label{sec:params}
%======================================================================

Table~\ref{tab:refine_params} lists all refinement-related
namelist parameters, including the new bypass parameters.

\begin{table}[h]
\centering
\caption{AMR refinement parameters in the \code{\&AMR\_PARAMS} and
\code{\&REFINE\_PARAMS} namelists.  New parameters are marked with $\star$.}
\label{tab:refine_params}
\small
\begin{tabular}{lll}
\toprule
Parameter & Default & Role \\
\midrule
\code{m\_refine}      & 8.0  & Lagrangian mass threshold (per level) \\
\code{sink\_refine}   & F    & Force max level around sinks \\
\code{err\_grad\_d}   & $-1$ & Density gradient threshold \\
\code{err\_grad\_p}   & $-1$ & Pressure gradient threshold \\
\code{err\_grad\_u}   & $-1$ & Velocity gradient threshold \\
\code{jeans\_refine}  & $-1$ & Jeans length in cells (per level) \\
\code{d\_jeans\_thre} & 0    & Density threshold for Jeans ($n_{\rm H}\,$cm$^{-3}$) \\
\code{r\_refine}      & $-1$ & Geometry region diameter (per level) \\
\code{a\_refine}      & 1.0  & Ellipticity $Y/X$ \\
\code{b\_refine}      & 1.0  & Ellipticity $Z/X$ \\
\code{exp\_refine}    & 2.0  & Norm exponent \\
\code{x/y/z\_refine}  & 0.5  & Region centre \\
\code{nexpand}        & 1    & Buffer zone width (per level) \\
\code{q\_refine\_holdback} & F & Cosmological holdback \\
$\star$ \code{jeans\_bypass\_holdback} & F & Enable Jeans bypass of holdback \\
$\star$ \code{n\_jeans\_bypass\_levels} & 1 & Max extra levels beyond $\ell_{\rm hb}$ \\
\bottomrule
\end{tabular}
\end{table}

The recommended configuration for cosmological runs with star formation is:
\begin{lstlisting}
&REFINE_PARAMS
m_refine          = 9*8.,
jeans_refine      = 9*4.,
jeans_bypass_holdback  = .true.
n_jeans_bypass_levels  = 1
d_jeans_thre      = 0.1         ! = n_star
q_refine_holdback = .true.
/
\end{lstlisting}


%======================================================================
\section{Test Results}
\label{sec:tests}
%======================================================================

All tests use the HR5 medium-box setup: $\ell_{\rm min}=9$,
$\ell_{\rm max}=17$, 12 MPI ranks, OMP=5, restarting from output 5
($a=0.164$, $z\approx 5.1$, nstep\_coarse=241).

\subsection{Bit-identical verification}

With \code{jeans\_bypass\_holdback=.false.} (default) and no
\code{jeans\_refine}, the new binary produces results identical to
the reference run: econs=$6.13\times10^{-4}$ at step 242,
econs=$6.71\times10^{-4}$ at step 243 (Job 122643).

\subsection{Unlimited bypass ($n_{\rm bypass}=\infty$, $\code{d\_jeans\_thre}=0$)}

With \code{jeans\_refine=9*4.},
\code{jeans\_bypass\_holdback=.true.}, and no level limit:
cascade refinement from level 14 through 17 within a single coarse
step.  Memory grew from 81\% to 88.5\% before OOM crash at fine step
1580.  Fine-step timestep decreased by $8\times$ due to subcycling
(Job 122644).

\subsection{Limited bypass ($n_{\rm bypass}=1$, $\code{d\_jeans\_thre}=0$)}

With \code{n\_jeans\_bypass\_levels=1}: only level 14 is bypassed,
cells at level 15 may be created but no further cascade.
Results: Job 122663 (pending).

\subsection{Future: $\code{d\_jeans\_thre}=\nstar$}

A test with \code{d\_jeans\_thre=0.1} will verify the physics-based
self-regulation, potentially allowing $n_{\rm bypass} > 1$ without
cascade, since only gas in the polytropic regime would trigger the
Jeans criterion.


%======================================================================
\section{Discussion}
\label{sec:discussion}
%======================================================================

\subsection{Physical motivation}

The holdback mechanism was designed to prevent over-refinement from the
mass criterion at early times.  However, it inadvertently blocks the
Jeans criterion, which has a different physical motivation: resolving
the Jeans length prevents artificial fragmentation and ensures
numerically converged star formation.

By allowing the Jeans criterion to operate ahead of the holdback
release, we expect:
\begin{itemize}
  \item Dense gas that will form stars is pre-refined before the
    holdback releases.
  \item When the holdback releases, the transition is smoother because
    the densest regions are already at the appropriate resolution.
  \item The SFR curve shows a gradual onset rather than an abrupt step.
\end{itemize}

\subsection{Interplay of the three parameters}

The bypass is controlled by three parameters that form a layered
defence:

\paragraph{Layer 1: \code{jeans\_bypass\_holdback} (on/off)}
Master switch.  When off (default), the system is identical to
standard \ramses.

\paragraph{Layer 2: \code{d\_jeans\_thre} (physics-based regulation)}
Restricts the Jeans check to dense gas where the polytropic floor
operates.  In this regime, $\lJ \propto \rho^{-1/6}$ varies slowly,
and 1--2 extra refinement levels suffice to resolve it.  Provides
self-consistent, physics-driven cascade termination.

\paragraph{Layer 3: \code{n\_jeans\_bypass\_levels} (hard safety limit)}
Absolute cap on the number of extra levels.  Prevents OOM regardless
of physics configuration.  Essential when
\code{d\_jeans\_thre=0} (all cells checked).


\bibliographystyle{plainnat}
\begin{thebibliography}{2}
\bibitem[Truelove et al.(1997)]{Truelove1997}
Truelove, J.~K., et al., 1997, ApJ, 489, L179

\bibitem[Dubois et al.(2021)]{Dubois2021}
Dubois, Y., et al., 2021, A\&A, 651, A109
\end{thebibliography}

\end{document}
